\SourceLeaf{86}%
\Article{13}
When we compare the examples in the preceding article, at least at first sight no simple criterion seems to offer itself for distinguishing the first moduli from the second. Nevertheless, there are \emph{two} such criteria, remarkable for their elegance and simplicity. The following considerations will prepare the way to one of them.

The modulus $p$, being a prime of the form $8n+1$, can be expressed in the form $aa+2bb$ in exactly one way (\emph{Disquiss. Arithm.}, Article 182, II). We assume that the roots $a,b$ are taken positive. Clearly $a$ is odd and $b$ even; let us put $b=2^\lambda c$, with $c$ odd. We now observe:

I. Since $p\equiv aa\ (\mathrm{mod.}\ c)$, $p$ is a quadratic residue of $c$, and hence also of each prime factor into which $c$ is resolved. Conversely, by the fundamental theorem, each of these prime factors is a quadratic residue of $p$, and so their product $c$ is also a quadratic residue of $p$. Since this also holds for $2$, $b$ is clearly a quadratic residue of $p$, and hence both $bb$ and $-bb$ are biquadratic residues.

II. Thus $-2bb$ must belong to the same class as $2$. Since $aa\equiv-2bb$, $2$ clearly belongs to class $A$ or $C$ according as $a$ is a quadratic residue of $p$ or a quadratic nonresidue.

III. Now suppose that $a$ has been resolved into prime factors. Denote those of the form $8m+1$ or $8m+7$ by $\alpha,\alpha',\alpha''$, etc., and those of the form $8m+3$ or $8m+5$ by $\beta,\beta',\beta''$, etc.; let the number of the latter be $=\mu$. Since $p\equiv2bb\ (\mathrm{mod.}\ a)$, $p$ is a quadratic residue of those prime factors of $a$ of which $2$ is a quadratic residue, that is, of $\alpha,\alpha',\alpha''$, etc.; it is a quadratic nonresidue of those factors of which $2$ is a quadratic nonresidue, that is, of $\beta,\beta',\beta''$, etc. Conversely, by the fundamental theorem, each of $\alpha,\alpha',\alpha''$, etc., is a quadratic residue of $p$, while each of $\beta,\beta',\beta''$, etc., is a quadratic nonresidue. We therefore conclude that their product $a$ is a quadratic residue of $p$ or a nonresidue according as $\mu$ is even or odd.

IV. It is easy to verify that the product of all $\alpha,\alpha',\alpha''$, etc., has the form $8m+1$ or $8m+7$, and that the same holds for the product of all $\beta,\beta',\beta''$, etc., when their number is even. In this case the product $a$ must therefore also have the form $8m+1$ or $8m+7$. In contrast, the product of all $\beta,\beta',\beta''$, etc.,\space%
\SourceLeaf{87}%
when their number is odd, has the form $8m+3$ or $8m+5$, and the same therefore holds for the product $a$ in this case.

From all this we obtain the elegant theorem:

\emph{Whenever $a$ has the form $8m+1$ or $8m+7$, $2$ belongs to $A$; whenever $a$ has the form $8m+3$ or $8m+5$, $2$ belongs to $C$.}

\Needspace{11\baselineskip}\begingroup\rightskip=0pt plus1fil\relax The examples listed in the preceding article confirm this. The first moduli decompose as follows: \mbox{$73=1+2\pdot36$}, \mbox{$89=81+2\pdot4$}, \mbox{$113=81+2\pdot16$}, \mbox{$233=225+2\pdot4$}, \mbox{$257=225+2\pdot16$}, \mbox{$281=81+2\pdot100$}, \mbox{$337=49+2\pdot144$}, \mbox{$353=225+2\pdot64$}; the second as follows: \mbox{$17=9+2\pdot4$}, \mbox{$41=9+2\pdot16$}, \mbox{$97=25+2\pdot36$}, \mbox{$137=9+2\pdot64$}, \mbox{$193=121+2\pdot36$}, \mbox{$241=169+2\pdot36$}, \mbox{$313=25+2\pdot144$}, \mbox{$401=9+2\pdot196$}, \mbox{$409=121+2\pdot144$}, \mbox{$433=361+2\pdot36$}, \mbox{$449=441+2\pdot4$}, \mbox{$457=169+2\pdot144$}.\par\endgroup

\Article{14}
The decomposition of $p$ into a square and twice a square has revealed such a remarkable connection with the classification of $2$ that it seems worthwhile to ask whether a decomposition into two squares, which $p$ is also known to admit, might give a similar result. Here are the decompositions of the numbers $p$ for which $2$ belongs to class
\[
\begin{array}{c|c}
A&C\\\hline
9+64&1+16\\
25+64&25+16\\
49+64&81+16\\
169+64&121+16\\
1+256&49+144\\
25+256&225+16\\
81+256&169+144\\
289+64&1+400\\
&9+400\\
&289+144\\
&49+400\\
&441+16
\end{array}
\]

\SourceLeaf{88}%
First we observe that, of the two squares into which $p$ decomposes, one must be odd, which we put $=aa$, and the other even, which we put $=bb$. Since $aa$ has the form $8n+1$, values of $b$ that are twice an odd number clearly correspond to values of $p$ of the form $8n+5$, which are excluded from our inductive investigation here, since they would have $2$ in class $B$ or $D$. For values of $p$ of the form $8n+1$, however, $b$ must be divisible by four. If we may trust the inductive evidence presented by this table, $2$ belongs to class $A$ for every modulus for which $b$ has the form $8n$, and to class $C$ for every modulus for which $b$ has the form $8n+4$. This theorem, however, requires much deeper investigation than the one obtained in the preceding article. Its proof must be preceded by several preliminary investigations concerning the order in which the numbers in $A,B,C,D$ follow one another.

\Needspace{22\baselineskip}\Article{15}
Let us denote the number of numbers in $A$ immediately followed by a number in $A,B,C,D$, respectively, by $(00),(01),(02),(03)$; likewise the number of numbers in $B$ followed by a number in $A,B,C,D$, respectively, by $(10),(11),(12),(13)$. Similarly, let $C$ contain respectively $(20),(21),(22),(23)$ numbers, and $D$ contain $(30),(31),(32),(33)$ numbers, followed by a number in $A,B,C,D$. We propose to determine these sixteen counts a priori. To help readers compare the general reasoning with examples, it seemed useful to give here the numerical values of the entries in the scheme $(S)$
\[
\begin{array}{llll}
(00),&(01),&(02),&(03)\\
(10),&(11),&(12),&(13)\\
(20).&(21),&(22),&(23)\\
(30),&(31),&(32),&(33)
\end{array}
\]
for each modulus whose classification we gave in Article 11.

\[
\begin{array}{c|c|c|c}
\begin{gathered}
p=5\\[3pt]
\begin{array}{l}
0,1,0\mathpunct{.}\,0\\
0,0,0,1\\
0,0,0,0\\
0,0,1,0
\end{array}
\end{gathered}&
\begin{gathered}
p=13\\[3pt]
\begin{array}{l}
0,1,2,0\\
1,1,0,1\\
0,1,0,1\\
1,0,1,1
\end{array}
\end{gathered}&
\begin{gathered}
p=17\\[3pt]
\begin{array}{l}
0,2,1,0\\
2,0,1,1\\
1,1,1,1\\
0,1,1,2
\end{array}
\end{gathered}&
\begin{gathered}
p=29\\[3pt]
\begin{array}{l}
2\mathpunct{.}\,3,0\mathpunct{.}\,2\\
1,1,2,3\\
2,1,2\mathpunct{.}\,1\\
1,2,3,1
\end{array}
\end{gathered}
\end{array}
\]

\SourceLeaf{89}%
\[
\begin{array}{c|c|c|c}
\begin{gathered}
p=37\\[3pt]
\begin{array}{l}
2,1,2,4\\
2,2,4,1\\
2,2,2,2\\
2,4,1,2
\end{array}
\end{gathered}&
\begin{gathered}
p=41\\[3pt]
\begin{array}{l}
0,4,3,2\\
4,2,2,2\\
3,2,3,2\\
2,2,2,4
\end{array}
\end{gathered}&
\begin{gathered}
p=53\\[3pt]
\begin{array}{l}
2,3,6,2\\
4,4,2,3\\
2,4,2,4\\
4,2,3,4
\end{array}
\end{gathered}&
\begin{gathered}
p=61\\[3pt]
\begin{array}{l}
4,3,2,6\\
3,3,6,3\\
4,3,4,3\\
3,6,3,3
\end{array}
\end{gathered}
\end{array}
\]

\[
\begin{array}{c|c|c}
\begin{gathered}
p=73\\[3pt]
\begin{array}{l}
5,6,4,2\\
6,2,5,5\\
4,5,4,5\\
2,5,5,6
\end{array}
\end{gathered}&
\begin{gathered}
p=89\\[3pt]
\begin{array}{l}
3,8,6,4\\
8,4,5,5\\
6,5,6,5\\
4,5,5,8
\end{array}
\end{gathered}&
\begin{gathered}
p=97\\[3pt]
\begin{array}{l}
2,6,7,8\\
6,8,5,5\\
7,5,7,5\\
8,5,5,6
\end{array}
\end{gathered}
\end{array}
\]

Since moduli of the forms $8n+1$ and $8n+5$ behave differently, we must treat them separately. We begin with the first.

\Article{16}
The symbol $(00)$ indicates how many different ways the equation $\alpha+1=\alpha'$ can be satisfied, with $\alpha,\alpha'$ denoting arbitrary numbers in $A$. For a modulus of the form $8n+1$, as understood here, $\alpha'$ and $p-\alpha'$ belong to the same collection. We may therefore say more elegantly that $(00)$ expresses the number of different ways of satisfying the equation $1+\alpha+\alpha'=p$. Clearly the congruence $1+\alpha+\alpha'\equiv0\ (\mathrm{mod.}\ p$ can also take the place of this equation.

Likewise,
\begin{center}
\begin{minipage}{\linewidth}
$(01)$ indicates the number of solutions of the congruence $1+\alpha+\beta\equiv0\ (\mathrm{mod.}\ p)$\par
$(02)$ the number of solutions of the congruence $1+\alpha+\gamma\equiv0$\par
$(03)$ the number of solutions of the congruence $1+\alpha+\delta\equiv0$\par
$(11)$ the number of solutions of the congruence $1+\beta+\beta'\equiv0$, etc.
\end{minipage}
\end{center}
where $\beta$ and $\beta'$ denote arbitrary numbers in $B$, $\gamma$ numbers in $C$, and $\delta$ numbers in $D$. We immediately obtain the following six equations:
\[
\begin{gathered}
(01)=(10),\quad(02)=(20),\quad(03)=(30),\\
(12)=(21),\quad(13)=(31),\quad(23)=(32)
\end{gathered}
\]

From any given solution of the congruence $1+\alpha+\beta\equiv0$ we obtain a solution of $1+\delta+\delta'\equiv0$ by taking for $\delta$ the number between $1\dotsfour p-1$\space%
\SourceLeaf{90}%
that satisfies $\beta\delta\equiv1$ (which clearly belongs to $D$), and for $\delta'$ the least positive residue of the product $\alpha\delta$ (which likewise belongs to $D$). The reverse passage from a given solution of $1+\delta+\delta'\equiv0$ to a solution of $1+\alpha+\beta\equiv0$ is equally clear: choose $\beta$ so that $\beta\delta\equiv1$, and at the same time put $\alpha\equiv\beta\delta'$. Thus the two congruences have the same number of solutions, or $(01)=(33)$.

Similarly, from the congruence $1+\alpha+\gamma\equiv0$ we derive $\gamma'+\gamma''+1\equiv0$ by choosing $\gamma'$ in $C$ so that $\gamma\gamma'\equiv1$, and $\gamma''$ in the same collection congruent to the product $\alpha\gamma'$. It follows easily that these two congruences admit the same number of solutions, or $(02)=(22)$.

Likewise, from $1+\alpha+\delta\equiv0$ we derive $\beta+\beta'+1\equiv0$ by choosing $\beta,\beta'$ so that $\beta\delta\equiv1$, $\beta\alpha\equiv\beta'$, giving $(03)=(11)$.

Finally, from $1+\beta+\gamma\equiv0$ we similarly derive both $\delta+1+\beta'\equiv0$ and $\gamma'+\delta'+1\equiv0$, and conclude that $(12)=(13)=(23)$.

We have thus obtained eleven equations among our sixteen unknowns, reducing them to five. The scheme $S$ can therefore be written as
\[
\begin{array}{llll}
h,&i,&k,&l\\
i,&l,&m,&m\\
k,&m,&k,&m\\
l,&m,&m,&i
\end{array}
\]

Three further equations of condition can easily be added. Every number in $A$, except the last, $p-1$, must be followed by a number in one of the collections $A,B,C$ or $D$. Thus
\[
(00)+(01)+(02)+(03)=2n-1
\]
and likewise
\[
\begin{aligned}
(10)+(11)+(12)+(13)&=2n\\
(20)+(21)+(22)+(23)&=2n\\
(30)+(31)+(32)+(33)&=2n
\end{aligned}
\]
In the notation just introduced, the first three equations give
\[
\begin{aligned}
h+i+k+l&=2n-1\\
i+l+2m&=2n\\
k+m&=n
\end{aligned}
\]

\SourceLeaf{91}%
The fourth is identical to the second. These equations allow us to eliminate three unknowns, so that all sixteen have now been reduced to two.

\Article{17}
To determine them completely, it will be useful to investigate the number of solutions of the congruence
\[
1+\alpha+\beta+\gamma\equiv0\ (\mathrm{mod.}\ p)
\]
where $\alpha,\beta,\gamma$ denote arbitrary numbers in $A,B,C$. Clearly $\alpha=p-1$ is inadmissible, since $\beta+\gamma\equiv0$ cannot hold. Substituting the remaining values for $\alpha$ in turn gives $h,i,k,l$ values of $1+\alpha$ belonging to $A,B,C,D$, respectively. For each \emph{given} value of $1+\alpha$ belonging to $A$, say $1+\alpha=\alpha^0$, the congruence $\alpha^0+\beta+\gamma\equiv0$ admits as many solutions as $1+\beta'+\gamma'\equiv0$ (put $\beta\equiv\alpha^0\beta'$, $\gamma\equiv\alpha^0\gamma'$), that is, $(12)=m$ solutions. Likewise, for each given value of $1+\alpha$ belonging to $B$, say $1+\alpha=\beta^0$, the congruence $\beta^0+\beta+\gamma\equiv0$ has as many solutions as $1+\alpha'+\beta'\equiv0$ (put $\beta\equiv\beta^0\alpha'$, $\gamma\equiv\beta^0\beta'$), that is, $(01)=i$ solutions. Similarly, for each given value of $1+\alpha$ belonging to $C$, say $1+\alpha=\gamma^0$, the congruence $\gamma^0+\beta+\gamma\equiv0$ can be solved in as many different ways as $1+\delta+\alpha'\equiv0$ (put $\beta\equiv\gamma^0\delta$, $\gamma\equiv\gamma^0\alpha'$), so the number of solutions is $(03)=l$. Finally, for each given value of $1+\alpha$ belonging to $D$, say $1+\alpha=\delta^0$, the congruence $\delta^0+\beta+\gamma\equiv0$ has as many solutions as $1+\gamma'+\delta'\equiv0$ (put $\beta\equiv\delta^0\gamma'$, $\gamma\equiv\delta^0\delta'$), that is, $(23)=m$ solutions. Adding all these counts, the congruence $1+\alpha+\beta+\gamma\equiv0$ clearly admits
\[
hm+ii+kl+lm
\]
different solutions.

In exactly the same way we find that, on substituting for $\beta$ the numbers in $B$ in turn, the sum $1+\beta$ takes respectively $(10),(11),(12),(13)$, or $i,l,m,m$, values belonging to $A,B,C,D$. For each \emph{given} value of $1+\beta$ belonging to these collections, the congruence $1+\beta+\alpha+\gamma\equiv0$ admits respectively $(02),(31),(20),(13)$, or $k,m,k,m$, different solutions. The total number of solutions is therefore
\SourceLeaf{92}%
\[
=ik+lm+km+mm
\]
We obtain the same value by basing the calculation on the values of the sum $1+\gamma$.

\Article{18}
These two expressions for the same count give the equation
\[
0=hm+ii+kl-ik-km-mm
\]
and hence, eliminating $h$ by means of $h=2m-k-1$,
\[
0=(k-m)^2+ii+kl-ik-kk-m
\]
The last two equations of Article 16 give $k=\frac12(l+i)$. Substituting this value turns $ii+kl-ik-kk$ into $\frac14(l-i)^2$, so the preceding equation, multiplied by $4$, becomes
\[
0=4(k-m)^2+(l-i)^2-4m
\]
Since $4m=2(k+m)-2(k-m)=2n-2(k-m)$, it follows that
\[
2n=4(k-m)^2+2(k-m)+(l-i)^2
\]
or
\[
8n+1=(4(k-m)+1)^2+4(l-i)^2
\]
Thus, putting
\[
4(k-m)+1=a,\quad2l-2i=b
\]
we have
\[
p=aa+bb
\]

It is known that $p$ can be decomposed into two squares in exactly one way, with the odd square taken as $aa$ and the even square as $bb$, so $aa,bb$ are completely determined numbers. But $a$ itself is also completely determined: the square root must be taken positive or negative according as the positive root has the form $4M+1$ or $4M+3$. We shall shortly discuss how to determine the sign of $b$.

Combining these new equations with the last three equations of Article 16 completely determines the five numbers $h,i,k,l,m$ in terms of $a,b$ and $n$, as follows:
\SourceLeaf{93}%
\[
\begin{aligned}
8h&=4n-3a-5\\
8i&=4n+a-2b-1\\
8k&=4n+a-1\\
8l&=4n+a+2b-1\\
8m&=4n-a+1
\end{aligned}
\]

If, in place of $n$, we prefer to use the modulus $p$, the scheme $S$, with every entry multiplied by $16$ to avoid fractions, is
\[
\begin{array}{l|l|l|l}
p-6a-11&p+2a-4b-3&p+2a-3&p+2a+4b-3\\
p+2a-4b-3&p+2a+4b-3&p-2a+1&p-2a+1\\
p+2a-3&p-2a+1&p+2a-3&p-2a+1\\
p+2a+4b-3&p-2a+1&p-2a+1&p+2a-4b-3
\end{array}
\]

\Article{19}
It remains to show how to choose the sign of $b$. In Article 10 we pointed out that the distinction between $B$ and $D$, which is not intrinsic, depends on the choice of $f$, for which either root of $xx\equiv-1$ must be taken, and that the two collections are interchanged when one root is replaced by the other. Inspection of the scheme just given shows that a similar interchange accompanies a change in the sign of $b$. We can therefore foresee that there must be a connection between the sign of $b$ and $f$. To find it, we first observe that, if $\mu$ denotes a nonnegative integer and $z$ takes all the values $1,2,3\dotsfour p-1$, then modulo $p$ we have either $\Sigma z^\mu\equiv0$ or $\Sigma z^\mu\equiv-1$, according as $\mu$ is not divisible by $p-1$ or is divisible by it. The latter part of the theorem follows from the fact that for a value of $\mu$ divisible by $p-1$, $z^\mu\equiv1$. We prove the former part as follows. Let $g$ denote a primitive root. All the values of $z$ coincide with the least residues of all $g^y$ when $y$ takes all the values $0,1,2,3\dotsfour p-2$, and hence $\Sigma z^\mu\equiv\Sigma g^{\mu\pdot y}$. Now
\[
\begin{gathered}
\Sigma g^{\mu y}=\frac{g^{\mu(p-1)}-1}{g^\mu-1},\quad\text{and hence}\\
(g^\mu-1)\Sigma z^\mu\equiv g^{\mu(p-1)}-1\equiv0
\end{gathered}
\]
For a value of $\mu$ not divisible by $p-1$, $g^\mu$ cannot be congruent to $1$, or equivalently $g^\mu-1$ cannot be divisible by $p$. Hence $\Sigma z^\mu\equiv0$. Q. E. D.

\newpage\SourceLeaf{94}%
Now expand the power $(z^4+1)^{\frac14(p-1)}$ by the binomial theorem. The preceding lemma gives
\[
\Sigma(z^4+1)^{\frac14(p-1)}\equiv-2\ (\mathrm{mod.}\ p)
\]
The least residues of all $z^4$ give all the numbers in $A$, each occurring four times. Among the least residues of $z^4+1$ we therefore have
\[
\begin{array}{rl}
4(00)&\text{to }A\\
4(01)&\text{to }B\\
4(02)&\text{to }C\\
4(03)&\text{to }D
\end{array}
\]
belonging to those classes, and four are $=0$ (namely, when $z^4\equiv p-1$). Using the criteria for the collections $A,B,C,D$, we obtain
\[
\begin{aligned}
\Sigma(z^4+1)^{\frac14(p-1)}\equiv{}&4(00)+4f(01)\\
&-4(02)-4f(03)
\end{aligned}
\]
and hence
\[
-2\equiv4(00)+4f(01)-4(02)-4f(03)
\]
or, substituting for $(00),(01)$, etc., the values found in the preceding article,
\[
-2\equiv-2a-2-2bf
\]
We therefore conclude that we must always have $a+bf\equiv0$, or, multiplying by $f$,
\[
b\equiv af
\]
This congruence determines the sign of $b$ when $f$ has already been chosen, or determines $f$ when the sign of $b$ is prescribed in some other way.
